% ══════════════════════════════════════════════════════════════════════════════
% Lecture 13: Graphs – DFS Theorems and Topological Sort
% ══════════════════════════════════════════════════════════════════════════════

% ─────────────────────────────────────────────────────────────────────────────
\subsection{Parenthesis Theorem}

DFS timestamps have a strict nested structure:

\begin{theorem}{Parenthesis Theorem}{parenthesis}
For any two vertices $u$ and $v$, exactly one of the following three situations occurs:
\begin{enumerate}
  \item The intervals $[u.d,\, u.f]$ and $[v.d,\, v.f]$ are \textbf{disjoint}, and
        neither $u$ nor $v$ is an ancestor of the other.
  \item $[v.d,\, v.f] \subset [u.d,\, u.f]$ (i.e., $u.d < v.d < v.f < u.f$)
        and $v$ is a \textbf{descendant} of $u$.
  \item $[u.d,\, u.f] \subset [v.d,\, v.f]$ (i.e., $v.d < u.d < u.f < v.f$)
        and $u$ is a \textbf{descendant} of $v$.
\end{enumerate}
\end{theorem}

\textbf{Illustration.} Using the example from Lecture 12 ($b$: 1/12, $a$: 2/7, $h$: 3/4,
$g$: 5/6, $c$: 8/11, $d$: 9/10, $e$: 13/16, $f$: 14/15):

\begin{center}
\begin{tikzpicture}[font=\small, >=Stealth]
  % Timeline
  \draw[->] (0,0) -- (17.5,0) node[right] {time};
  \foreach \t in {1,...,16}{
    \draw (\t, 0.08) -- (\t, -0.08);
    \node[below, font=\tiny] at (\t, -0.1) {\t};
  }
  % Intervals
  \foreach \label/\s/\e/\y/\col in {
    $b$/1/12/1.8/blue!30,
    $a$/2/7/1.2/green!30,
    $h$/3/4/0.6/orange!30,
    $g$/5/6/0.6/orange!30,
    $c$/8/11/1.2/green!30,
    $d$/9/10/0.6/orange!30,
    $e$/13/16/1.2/blue!20,
    $f$/14/15/0.6/orange!30}{
    \draw[fill=\col, draw=gray, rounded corners=1pt]
          (\s, \y-0.25) rectangle (\e, \y+0.25);
    \node[font=\tiny\bfseries] at ({(\s+\e)/2}, \y) {\label};
  }
\end{tikzpicture}
\end{center}

We see that the intervals are either nested (ancestor/descendant) or disjoint
(no descendant relationship). They never partially overlap.

% ─────────────────────────────────────────────────────────────────────────────
\subsection{White-Path Theorem}

\begin{theorem}{White-Path Theorem}{whitepath}
In a DFS forest of a graph $G$, vertex $v$ is a \textbf{descendant} of $u$
if and only if at time $u.d$ (when $u$ is just discovered), there exists a
path from $u$ to $v$ consisting \textbf{entirely of WHITE vertices}
(not yet discovered), with the exception of $u$ itself which has just been colored GRAY.
\end{theorem}

This theorem fully characterizes the descendant relationship in the DFS forest:
$v$ will be in $u$'s subtree precisely if such a white path exists
at the moment $u$ is discovered.

% ══════════════════════════════════════════════════════════════════════════════
\section*{Topological Sort}
\addcontentsline{toc}{subsection}{Topological Sort}
% ══════════════════════════════════════════════════════════════════════════════

% ─────────────────────────────────────────────────────────────────────────────
\subsection{Definition}

\begin{definition}{Topological Sort}{topo}
\textbf{Input:} a directed acyclic graph (DAG) $G = (V, E)$.

\textbf{Output:} a linear ordering of vertices such that for every edge
$(u, v) \in E$, $u$ appears before $v$ in the ordering.
\end{definition}

\textbf{Motivating example:} the order in which one gets dressed in the morning.
If you must put on socks before shoes, and underwear
before pants, a topological sort provides an order consistent with all these
constraints.

\begin{remark}{DAG and Applicability}{dagrem}
A topological sort exists only if the graph is \textbf{acyclic} (DAG).
In the presence of a cycle, no linear ordering can satisfy all
precedence constraints.
\end{remark}

% ─────────────────────────────────────────────────────────────────────────────
\subsection{Characterizing DAGs via DFS}

\begin{theorem}{DAG $\Leftrightarrow$ No DFS Back Edge}{daglem}
A directed graph $G$ is acyclic if and only if a DFS of $G$ produces
\textbf{no back edges}.
\end{theorem}

\begin{proof}
\textbf{Back edge $\Rightarrow$ cycle.}
Suppose a back edge $(u, v)$ exists. Then $v$ is an ancestor of $u$
in the DFS forest, so there is a path from $v$ to $u$ in $G$. The back edge
$(u, v)$ closes this path into a cycle.

\begin{center}
\begin{tikzpicture}[font=\small, >=Stealth,
    node/.style={circle, draw, fill=blue!10, minimum size=0.6cm, inner sep=1pt}]
  \node[node] (v) at (1, 2) {$v$};
  \node[node] (m) at (1, 1) {$\cdots$};
  \node[node] (u) at (1, 0) {$u$};
  \draw[->, keygreen, thick] (v) -- (m) node[midway, right, font=\tiny]{$T$};
  \draw[->, keygreen, thick] (m) -- (u) node[midway, right, font=\tiny]{$T$};
  \draw[->, keyred, thick, bend right=50] (u) to node[left, font=\tiny]{$B$} (v);
  \node[keyred, font=\tiny] at (2.5, 1) {cycle!};
\end{tikzpicture}
\end{center}

\textbf{Cycle $\Rightarrow$ back edge.}
Let $C$ be a cycle in $G$ and $v$ be the first vertex of $C$ discovered by DFS.
Let $(u, v)$ be the edge preceding $v$ in the cycle. At time $v.d$, the vertices
of $C$ form a white path from $v$ to $u$. By the white-path theorem,
$u$ becomes a descendant of $v$, so the edge $(u, v)$ is a back edge.
\end{proof}

% ─────────────────────────────────────────────────────────────────────────────
\subsection{Algorithm}

\begin{algorithm}[H]
\caption{Topological-Sort($G$)}
\begin{algorithmic}[1]
\State Call \textsc{DFS}($G$) to compute finishing times $v.f$ for all $v$
\State \Return the vertices in \textbf{decreasing} order of finishing times $v.f$
\end{algorithmic}
\end{algorithm}

\textbf{Efficient implementation:} explicit sorting is not necessary.
One can insert each vertex at the \emph{head} of a linked list at the moment
it is finished (colored BLACK). At the end of the DFS, the list contains the vertices
in topological order.

\begin{complexity}{Topological-Sort}{topocplx}
\textbf{Complexity: $\Theta(V + E)$} — identical to DFS complexity.
\end{complexity}

% ─────────────────────────────────────────────────────────────────────────────
\subsection{Proof of Correctness}

We must show that if $(u, v) \in E$, then $v.f < u.f$ (so $u$ appears before
$v$ in the decreasing order of finishing times).

When edge $(u,v)$ is explored during DFS, $u$ is GRAY. Three cases for $v$:

\begin{center}
\renewcommand{\arraystretch}{1.4}
\begin{tabular}{cll}
\toprule
\textbf{Color of $v$} & \textbf{Consequence} & \textbf{Result} \\
\midrule
GRAY & $v$ would be an ancestor of $u$ & Back edge $\Rightarrow$ cycle $\Rightarrow$ contradiction \\
WHITE & $v$ becomes descendant of $u$ & By parenthesis: $u.d < v.d < v.f < u.f$ \\
BLACK & $v$ already finished, $u$ not yet & Directly $v.f < u.f$ \\
\bottomrule
\end{tabular}
\end{center}

In all non-contradictory cases, we indeed obtain $v.f < u.f$. \qed

% ─────────────────────────────────────────────────────────────────────────────
\subsection{Example}

\begin{center}
\begin{tikzpicture}[font=\small, >=Stealth,
    node/.style={circle, draw, fill=blue!10, minimum size=0.75cm, inner sep=1pt}]
  % Nodes with DFS timestamps
  \node[node] (undies) at (0, 2)    {\small undies}; \node[above=2pt] at (undies.north) {\tiny 1/8};
  \node[node] (pants)  at (2, 2)    {\small pants}; \node[above=2pt] at (pants.north)  {\tiny 2/7};
  \node[node] (belt)   at (4, 2)    {\small belt}; \node[above=2pt] at (belt.north)   {\tiny 3/6};
  \node[node] (shirt)  at (0, 0)    {\small shirt}; \node[above=2pt] at (shirt.north)  {\tiny 9/12};
  \node[node] (tie)    at (2, 0)    {\small tie};  \node[above=2pt] at (tie.north)    {\tiny 10/11};
  \node[node] (jacket) at (4, 0)    {\small jacket}; \node[above=2pt] at (jacket.north) {\tiny 13/14};
  \node[node] (socks)  at (6, 1)    {\small socks}; \node[above=2pt] at (socks.north) {\tiny 15/16};
  \node[node] (shoes)  at (5, -1)   {\small shoes}; \node[above=2pt] at (shoes.north)  {\tiny 4/5};
  \draw[->,thick] (undies)--(pants);
  \draw[->,thick] (undies)--(shoes);
  \draw[->,thick] (pants)--(belt);
  \draw[->,thick] (pants)--(shoes);
  \draw[->,thick] (belt)--(jacket);
  \draw[->,thick] (shirt)--(belt);
  \draw[->,thick] (shirt)--(tie);
  \draw[->,thick] (tie)--(jacket);
  \draw[->,thick] (socks)--(shoes);
  % Topological order arrow
  \node[draw=none, fill=none, keyred, font=\bfseries\small] at (3, -2)
        {Topo. Order: $f=16,15,14,12,11,8,7,6,5$};
\end{tikzpicture}
\end{center}

The topological sort (decreasing finishing times) yields a valid
dressing order: each garment appears after those that must be put on before it.

% ─────────────────────────────────────────────────────────────────────────────
\subsection{Lecture Summary}

\begin{tcolorbox}[colback=lightblue, colframe=keyblue, fonttitle=\bfseries,
                  title={Key Points}]
\begin{itemize}[noitemsep]
  \item \textbf{Parenthesis Theorem:} the intervals $[u.d, u.f]$ are either
        nested or disjoint — never partially overlapping.
  \item \textbf{White-Path Theorem:} $v$ is a descendant of $u$ iff there exists
        a white path from $u$ to $v$ at time $u.d$.
  \item \textbf{DAG $\Leftrightarrow$ No back edge:} a tool for detecting
        cycles.
  \item \textbf{Topological Sort:} DFS + decreasing order of finishing times.
        Complexity $\Theta(V+E)$.
\end{itemize}
\end{tcolorbox}